\section{Introduction}
\label{sec:introduction}

Active causal discovery is a sequential inference problem: an agent must extract tentative structure from observational data, choose interventions under a finite budget, integrate the resulting evidence, and finally commit to a directed acyclic graph (DAG). Unlike purely observational causal discovery, the target is not fully determined by the observed distribution alone. In general, observational data identify a Markov equivalence class rather than a unique DAG, so additional interventions are required to resolve the remaining directional ambiguity \citep{spirtes2000causation,verma1990equivalence,andersson1997characterization,eberhardt2005interventions}. This makes active causal discovery a natural agent task rather than a static prediction problem.

Structural causal models (SCMs) make this domain especially well suited to diagnostic evaluation. They provide an exact hidden causal state together with precise intervention semantics, while CPDAG and equivalence-class theory make the observational ceiling explicit \citep{spirtes2000causation,andersson1997characterization,hauser2012gies}. The resulting separation is unusually clean: a system can fail by missing adjacencies, by failing to recover directions that are already compelled observationally, or by failing to use interventions to orient the remaining edges. In other words, SCM theory gives the evaluator access to distinct theoretical objects against which different parts of the task can be scored separately.

Most contemporary agent benchmarks instead report a single end-to-end success signal on complex environments such as software engineering, web navigation, and computer use \citep{jimenez2024swebench,liu2024agentbench,zhou2024webarena,xie2024osworld,mialon2023gaia}. Those evaluations are valuable, but they are less diagnostic when the task itself contains separable inference stages. In causal discovery specifically, prior work has examined whether language models reason about causality at all \citep{kiciman2023llmcausal,zecevic2023causalparrots}, but a single graph score on the active task can still confound observational structure recovery, interventional reasoning, and final graph commitment. Benchmark conclusions can also become hard to interpret when graph families induce a non-trivial blind-random score floor. What is missing is not another leaderboard, but an evaluation instrument whose scoring layers and calibration are aligned with the underlying causal objects.

We introduce the \textbf{Active Causal Discovery Benchmark (ACDB)}, an SCM-grounded benchmark that fixes a hidden linear--Gaussian world generator, an observe--intervene--submit interaction protocol, and a paired seed map across agents. ACDB evaluates submissions with a layered scoring contract over skeleton recovery, observationally compelled orientation, full directed recovery, and structural error, and it calibrates its ladder with level-specific expected/random directed-F1 floors. We use this benchmark to compare classical baselines and LLM-based policies under the same active protocol, including end-to-end raw-data agents and hybrid policies that isolate observational structure recovery from interventional orientation.

On the current five-model, five-ablation panel, ACDB reveals several separable patterns. Layered scores expose distinctions that a single final-graph metric would hide; supplying a strong observational prior often yields large gains but does not remove the interventional orientation burden; and granting conditional-independence-test tools does not reliably improve directed recovery. More broadly, performance varies substantially across model families and degrades asymmetrically with graph complexity relative to the classical active baseline. These findings are not the benchmark's primary contribution. Their role in this paper is to demonstrate that ACDB can localize failure to observational recovery, interventional orientation, or end-to-end integration.

Our contributions are:
\begin{itemize}
    \item We introduce ACDB, a benchmark for active causal discovery with a fixed SCM generator, a fixed observe--intervene--submit protocol, and a theory-grounded layered scoring contract.
    \item We present a calibration methodology for active causal-discovery evaluation based on expected/random directed-F1 floors and a graph-scale ladder designed to avoid structurally guessable regions.
    \item We provide an empirical study across classical and LLM-based policies showing that the benchmark can localize failure across observational recovery, interventional orientation, and end-to-end graph integration.
\end{itemize}
